\SourceLeaf{14}
\begin{center}\small\setlength{\tabcolsep}{3pt}\begin{tabular}{@{}rll@{}}
10)&$m=-7$&$a,\qquad b$ per 7\\
11)&$m=+10$ si&$2a+5b,\qquad b,\qquad a\pm7b$ per 20\\
12)&$m=-10$&$2a+5b,\qquad b,\qquad a\pm3b$ per 20\\
13)&$m=+11$&$2a+11b,\qquad b,\qquad2a\pm3b,\qquad a\pm b$ per 22\\
14)&$m=-11$&$b,\qquad a\pm4b$ per 11.
\end{tabular}\end{center}
\begin{center}[3.]\par\smallskip Die cubischen Reste.\end{center}
\noindent\hspace*{\parindent}$p$ eine Primzahl der Form $3m+1$
\[4p=aa+27bb,\]
\par\noindent\hspace*{\parindent}$q$ eine andere Primzahl der Form $3n\pm1$
\[(a+3b\sqrt{-3})^{\frac{q\mp1}{3}}\equiv A\pm B\sqrt{-3}.\]
\par\noindent\hspace*{\parindent}Nun gibt es drei Fälle
\[\begin{array}{rlll}
\text{I.}&B\equiv0\quad\;(\mathrm{mod.}\ q).&\text{Dann ist }q^{\frac{p-1}{3}}\equiv1&\;(\mathrm{mod.}\ p).\\[5pt]
\text{II.}&B\equiv A\quad\;(\mathrm{mod.}\ q).&\text{Dann ist }q^{\frac{p-1}{3}}\equiv-\tfrac12-\tfrac12\dfrac{a}{3b}&\;(\mathrm{mod.}\ p).\\[5pt]
\text{III.}&B\equiv-A\quad\;(\mathrm{mod.}\ q).&\text{Dann ist }q^{\frac{p-1}{3}}\equiv-\tfrac12+\tfrac12\dfrac{a}{3b}&\;(\mathrm{mod.}\ p).
\end{array}\]
\Needspace{9\baselineskip}
\noindent\hspace*{\parindent}Erstes Beispiel.\qquad $\begin{array}[t]{ll}p=13,&q=5\\a=5,&b=1\end{array}$
\[(5+3\sqrt{-3})^2=-2+30\sqrt{-3},\quad B\equiv0\ \mathrm{mod.}\ 5\ \text{also Casus I}\]
\[5^4=625\equiv1.\]
\Needspace{9\baselineskip}
\noindent\hspace*{\parindent}Zweites Beispiel.\qquad $\begin{array}[t]{ll}p=19,&q=5\\a=7,&b=1\end{array}$
\[(7+3\sqrt{-3})^2=22+42\sqrt{-3},\quad B\equiv-A\ \text{also Casus III}\]
\[5^6\equiv7\quad\mathrm{mod.}\ 19,\qquad42\equiv-3+7\quad\mathrm{mod.}\ 19.\]
